\documentclass[11pt,a4paper]{article}
\usepackage[T1]{fontenc}
\usepackage{lmodern}
\usepackage[margin=26mm]{geometry}
\usepackage{amsmath,amssymb,amsthm}
\usepackage{microtype}
\usepackage{enumitem}
\usepackage[hidelinks]{hyperref}
\setlength{\parskip}{3pt}
\setlength{\emergencystretch}{3em}
\newtheorem{theorem}{Theorem}[section]
\newtheorem{lemma}[theorem]{Lemma}
\newtheorem{corollary}[theorem]{Corollary}
\theoremstyle{remark}
\newtheorem{remark}[theorem]{Remark}
\newcommand{\Z}{\mathbb Z}
\newcommand{\N}{\mathbb N}
\newcommand{\Q}{\mathbb Q}
\newcommand{\F}{\mathbb F}
\newcommand{\ps}[1]{#1[[x]]}
\hypersetup{pdftitle={A389475: Two cubic congruences in characteristic three},pdfauthor={Dev Goyal}}
\title{A389475\\[3pt]Two cubic congruences in characteristic three}
\author{Dev Goyal}
\date{24 October 2026}
\begin{document}
\maketitle
\begin{center}
\fbox{\begin{minipage}{0.91\linewidth}\small
\textbf{AI contribution and review status.}
ChatGPT (OpenAI) generated the mathematical exposition, proof development,
Lean interfaces, and checks for this collection in response to Dev Goyal's
prompts. Dev Goyal is the sole human author. Independent human review is
pending. Original sequence authors, known mathematical results, and adapted
upstream code retain their separate credits. This document makes no claim
of publication, OEIS acceptance, or worldwide priority.
\end{minipage}}
\end{center}
\begin{abstract}
For the unique normalized integer solution of $V(3x^2)=2x^2+V^2-2V^3$, we prove both recorded modulus-three identities. The construction is integral and unique; the congruence argument is a reduction of the source equation followed by standard Frobenius.
\end{abstract}
\section{Source and conclusions}
The A389475 series \cite{oeis} is the normalized integer solution
\[
V(3x^2)=2x^2+V(x)^2-2V(x)^3,\qquad V=x+O(x^2).
\]
The offset is one. We prove the two coefficientwise identities
\[
V(x)^3\equiv V(x^3)\pmod3,\qquad
V(x)^2\equiv x^2+2V(x^3)\pmod3.
\]
The first is standard Frobenius; the second uses this source equation.
\section{Integral construction and uniqueness}

\begin{theorem}\label{thm:cubicexist}
For every $s\in\Z$, there is a unique normalized series
$R_s\in\ps{\Z}$ satisfying
\begin{equation}\label{eq:cubicfamily}
R_s(3x^2)=2x^2+R_s(x)^2+2sR_s(x)^3.
\end{equation}
\end{theorem}

\begin{proof}
For any normalized integral series $U$, define its residual by
\[
E_s(U)=U^2+2sU^3-U(3x^2)+2x^2.
\]
Every coefficient of $E_s(U)$ is even. Indeed, after reduction modulo
two, Frobenius gives $U(x)^2=U(x^2)$ and $3x^2=x^2$, so the residual
vanishes. Also $E_s(U)=O(x^3)$ for every normalized $U$.

Suppose $V-U=cx^r+O(x^{r+1})$, with $r\ge2$. The quadratic difference
has leading term $2cx^{r+1}$; the cubic difference has order at least
$r+2$; and $(V-U)(3x^2)$ has order at least $2r\ge r+2$.
Consequently
\begin{equation}\label{eq:triangular}
E_s(V)-E_s(U)=2cx^{r+1}+O(x^{r+2}).
\end{equation}
Starting with $U=x$, choose its degree-$r$ coefficient to cancel the
degree-$(r+1)$ coefficient of the residual. The needed division by two
is exact by the preceding parity argument. Equation~\eqref{eq:triangular}
shows that earlier cancellations remain valid. Each degree is thus
determined by finitely many earlier coefficients, and the resulting
integral series has zero residual in every degree.

For uniqueness, take the first degree $r\ge2$ where two candidate
solutions differ. Equation~\eqref{eq:triangular} would give a nonzero
degree-$(r+1)$ difference of residuals, because $2c\ne0$ in $\Z$.
This contradicts both residuals being zero.
\end{proof}

The proof formalizes divisibility of the residual by two before using
integer division. Thus the construction does not assume the existence
or integrality of an OEIS sequence. The names used below denote the
uniquely constructed series.


Taking $s=-1$ in Theorem~\ref{thm:cubicexist} constructs the desired
series $V=R_{-1}$ and proves uniqueness for its literal source equation.
\section{The standard Frobenius identity}
For any integer formal series $W$ and prime $p$, coefficient reduction gives
\[
\overline W(x)^p=\overline W(x^p)\qquad\text{in }\F_p[[x]].
\]
For each fixed coefficient one may truncate $W$ at that degree. The
polynomial identity $(u+v)^p=u^p+v^p$ in characteristic $p$, together
with $a^p=a$ in $\F_p$, proves the displayed identity coefficient by
coefficient. Thus it holds for every formal series, without a convergence
assumption. This is standard algebra; its application below is not a new
general principle.

\section{The two modulus-three identities}
\begin{theorem}
For every coefficient index $n\ge0$,
\begin{align*}
[x^n]\bigl(V(x)^3-V(x^3)\bigr)&\equiv0\pmod3,\\
[x^n]\bigl(V(x)^2-x^2-2V(x^3)\bigr)&\equiv0\pmod3.
\end{align*}
\end{theorem}
\begin{proof}
Apply Frobenius with $p=3$ and $W=V$ for the first identity.
For the second reduce the source equation coefficientwise to $\F_3$.
Since $V(0)=0$, its left side reduces to $V(0)=0$. Hence
\[
0=2x^2+V^2-2V^3,\qquad V^2=x^2+2V^3.
\]
Substitute the first identity to obtain $V^2=x^2+2V(x^3)$.
Equality of series gives every stated coefficient congruence.
\end{proof}
This proves the assertions labeled (4.a) and (4.b) in the retrieved
entry. Both hold for the full series, including degree zero.
\section{Lean proof and exact scope}
The dedicated entry file is \path{lean/Entries/A389475.lean}. Its declarations lie in namespace \texttt{FinalCollection.A389475}. The principal checked endpoints are:
\begin{itemize}[nosep]
\item \texttt{source\_equation}: the literal defining equation.
\item \texttt{source\_existence\_uniqueness}: the integral source construction and uniqueness.
\item \texttt{integral\_family\_existence}: the integer-parameter family used in the construction.
\item \texttt{cube\_mod\_three}: the full Frobenius identity.
\item \texttt{square\_mod\_three}: the full square identity derived from the source.
\item \texttt{cube\_coefficients}: the first assertion at every coefficient.
\item \texttt{square\_coefficients}: the second assertion at every coefficient.
\end{itemize}
The full implementations are in \path{lean/More/CubicConsequences.lean}.
Every index is universally quantified. The first congruence is a standard identity valid for every integral formal power series; the second is a consequence of this particular source equation. The companion also separately states normalization and source uniqueness.

The companion is a separate, explicitly typed Lean entry file. Its proofs
invoke the preserved implementation modules; the complete dependency source
is included in the same archive. Shared lemmas are reused rather than
assumed. The entry file is not a self-contained replacement for Mathlib or
the supplied common library.

The project pins Lean 4.19.0 and Mathlib revision
\begin{center}\small\ttfamily
c44e0c8ee63ca166450922a373c7409c5d26b00b.
\end{center}
The collection checker recompiles the proof sources and audits the
theorems' transitive axiom dependencies. Only \texttt{propext},
\texttt{Classical.choice}, and \texttt{Quot.sound} are permitted.
The delivered verification reports record the actual successful runs.
Finite numerical checks are corroboration; the displayed conclusions are
unbounded theorems. An AI review of source correspondence does not replace
independent human review.
\begin{thebibliography}{9}
\raggedright
\bibitem{oeis} OEIS Foundation, entry A389475,
\url{https://oeis.org/A389475}. Source checked 24 October 2026.
The source sequence and selected conjectures are credited to Paul D. Hanna.
\bibitem{repo} D. Goyal, supplied reference repository,
\url{https://github.com/dvgyl/oeis-a396843-proof}.
This separate presentation belongs to the final collected proof package;
no earlier DOI is reassigned to it.
\end{thebibliography}
\end{document}
